% 16. Źródła i licencje
\section{Źródła i licencje}\label{sec:zrodla}

\begin{instrukcja}[wymagane]
Skąd pochodzą zdjęcia, ikony, fonty, biblioteki i gotowe fragmenty kodu, na jakiej licencji. Jeśli
wszystko jest własne - napiszcie to. Literaturę podajcie w formie listy.
\end{instrukcja}

\subsection{Materiały i biblioteki}
\begin{table}[H]\centering
\caption{Wykorzystane materiały i ich licencje}\label{tab:licencje}
{\small
\begin{tabularx}{\linewidth}{@{}p{3cm}p{3.4cm}p{2.6cm}L@{}}
\naglowektabeli \nt{Zasób} & \nt{Nazwa} & \nt{Licencja} & \nt{Źródło} \\
ikony & Lucide (\texttt{lucide-react}) & ISC & \url{https://lucide.dev} \\
font & systemowy (Segoe UI, Roboto) & -- & nie pobieramy fontów \\
zdjęcia & brak & -- & -- \\
efekty & wzory: Magic UI (Number Ticker, Border Beam), React Bits (Spotlight Card) & MIT & \url{https://magicui.design}, \url{https://reactbits.dev}; kod napisany od nowa \\
biblioteki & Astro, React, motion & MIT & \url{https://astro.build}, \url{https://react.dev}, \url{https://motion.dev} \\
\bottomrule
\end{tabularx}}
\end{table}

\subsection{Literatura}
\begin{enumerate}[label={[\arabic*]}]
  \item J. Tidwell, \emph{Projektowanie interfejsów. Sprawdzone wzorce projektowe}, wyd. II, Helion, 2020.
  \item Prezentacje z wykładów 1-4, Metody opracowywania interfejsów użytkownika, UKEN 2026/27.
  \item Dokumentacja Astro, \url{https://docs.astro.build}; dokumentacja React, \url{https://react.dev}.
\end{enumerate}
